\section{Вводная часть}                     % Раздел 1
Вспомним некоторые утверждения, частично доказанные в прошлом семестре:
\begin{proposition}
    Набор строк матрицы $A$ линейно зависим $\Leftrightarrow$ существует строка $\vec{c}\neq 0$, такая что $cA=0$.
\end{proposition}
\begin{proof}
    Пусть $n$ - порядок матрицы $A$.
    \begin{itemize}
        \item[$\Leftarrow$] Пусть $\vec{c}=\begin{pmatrix} \alpha_1, \dots, \alpha_n\end{pmatrix}$. Тогда $\vec{c}A=\sum_{i=1}^{n} \alpha_i \vec{a_i}=0$, где $\vec{a_i}$ - строки матрицы $A$. Следовательно, строки матрицы $A$ линейно зависимы.
        \item[$\Rightarrow$] По определению $\exists$ $\alpha_1, \dots, \alpha_n:\sum_{i=1}^{n} \alpha_i^2 >0 \hookrightarrow \sum_{i=1}^{n} \alpha_i \vec{a_i}=0$, где $\vec{a_i}$ - строки матрицы $A$. Значит $\vec{c}=\begin{pmatrix} \alpha_1, \dots, \alpha_n\end{pmatrix}$ удовлетворяет условию утверждения.
    \end{itemize}
\end{proof}
\begin{proposition}
    Пусть $A, B$ - квадратные матрицы порядка $n$. Если у $A$ набор строк линейно зависим, то у $AB$ набор строк тоже линейно зависим.
\end{proposition}
\begin{proof}
    По предыдущему утверждению $\exists \ \vec{c}\hookrightarrow \vec{c}A=0$. Тогда $(\vec{c}A)B=\vec{c}(AB)=0$, а значит набор строк $AB$ линейно зависим. \textit{Замечание: отметим, что можно брать одинаковые коэффициенты в нетривиальные линейные комбинации для $A$ и для $AB$}
\end{proof}
\begin{proposition}
    Элементарные преобразования строк переводят линейно зависимые наборы строк в линейно зависимые, а линейно независимые в линейно независимые.
\end{proposition}
\begin{proof}
    Пусть дан набор $a_1, \dots, a_n$. 
    \begin{itemize}[leftmargin=*]
        \item Докажем, что ЛНЗ наборы переходят в ЛНЗ наборы: \begin{itemize}
            \item Пусть ЛНЗ набор $a_1, \dots, a_n \rightarrow$ $a_1+a_2, a_2, \dots, a_n$ - ЛЗ набор, тогда $\alpha_1 (a_1+a_2) + \alpha_2 a_2 + \dots +\alpha_n a_n$ = 0. Но из линейной независимости $a_1, \dots, a_n$: $\alpha_1 = \alpha_1 + \alpha_2 = \alpha_3 = \dots =\alpha_n = 0$. Тогда  $\alpha_1 (a_1+a_2) + \alpha_2 a_2 + \dots +\alpha_n a_n$ - тривиальная линейная комбинация, противоречие.
            \item Пусть ЛНЗ набор $a_1, \dots, a_n \rightarrow$ $\lambda a_1, a_2, \dots, a_n$ - ЛЗ набор, тогда $\alpha_1 \lambda a_1 + \alpha_2 a_2 + \dots +\alpha_n a_n$ = 0. Но из линейной независимости $a_1, \dots, a_n$: $\alpha_1 \lambda = \alpha_2 = \dots =\alpha_n = 0$. Тогда $\alpha_1 \lambda a_1 + \alpha_2 a_2 + \dots +\alpha_n a_n$ - тривиальная линейная комбинация, противоречие.
        \end{itemize}
        \item Так как ЛНЗ наборы переходят в ЛНЗ и элементарные операции обратимы, то ЛЗ наборы переходят в ЛЗ.
    \end{itemize}
\end{proof}
\begin{proposition}
    При элементарных операциях со столбцами сохраняются коэффициенты линейных зависимостей между строками, а линейно независимые строки остаются линейно независимыми.
\end{proposition}
\begin{proof}
    \begin{itemize}
        \item Пусть строки матрицы $A$ ЛЗ, тогда $\exists \ \vec{c} \hookrightarrow \vec{c}A=0$. Умножая на элементарную матрицу $S$ справа (операция со столбцами) получаем $\vec{c}AS=(\vec{c}A)S=0$. Следовательно, строки $AS$ - ЛЗ.
        \item Пусть строки матрицы $A$ ЛНЗ, тогда если после элементарного преобразования столбцов строки матрица $AS$ cтали ЛЗ, то при обратном элементарном преобразовании ЛЗ строки матрицы $AS$ переходят в ЛНЗ, что противоречит предыдущему пункту.
    \end{itemize}
\end{proof}
\begin{note}
    Аналогичное утверждение верно и для элементарных операций со строками и ЛЗ/ЛНЗ столбцов, так как операция транспонирования сохраняет ЛЗ/ЛНЗ.
\end{note}
\begin{theorem} [Крамера]
    Матрица $A_{n \times n}$ - невырожденная $\Rightarrow \forall$ столбец $d$ высоты $n$ может быть представлен в виде линейной комбинации столбцов матрицы $A$, причем единственным образом.
\end{theorem}
\begin{note}
    Аналогичное утверждение верно и для строк.
\end{note}
\section{Ранг матрицы}
\label{Ранг матрицы.}
\begin{definition}
    Пусть в матрице $A$ существует ЛНЗ набор из $r$ строк и $\forall k>r$ нельзя выбрать ЛНЗ набор из $k$ строк. Тогда $r$ называется строчным рангом. Обозначение $r_1 (A)$. У нулевой матрицы строчный ранг $= 0$. 
\end{definition}
\begin{note}
    Аналогично определяется $r_2 (A)$ - столбцовый ранг.
\end{note}
\begin{definition}
    Пусть дана матрица $A$. Рассмотрим ЛНЗ набор из $r_1 (A)$ строк. Будем называть его базисным, а его строки базисными.
\end{definition}
\begin{theorem}
    Любая строка матрицы есть линейная комбинация её базисных строк.
\end{theorem}
\begin{proof}
    Рассмотрим два случая:
    \begin{itemize}
        \item Если эта строка одна из базисных, то берем её линейную комбинацию с коэффициентом 1, а все остальные базисные строки с коэффициентами 0.
        \item Если эта строка $a$ небазисная:  пусть $a_1, \dots, a_r$ базисные строки, тогда $a_1, \dots, a_r, a$ - ЛЗ $\Rightarrow \exists\ \alpha_1 , \dots, \alpha_r, \alpha \hookrightarrow \alpha_1 a_1 + \dots + \alpha_r a_r + \alpha a = 0$. Тогда если $\alpha = 0$, то в силу нетривиальности линейной комбинации и $\alpha_1 a_1 + \dots + \alpha_r a_r = 0$ получаем, что базисный набор ЛЗ, противоречие. Значит $\alpha \neq 0$, следовательно $a=\sum_{i=1}^{r}-\frac{\alpha_i}{\alpha} a_i$.
    \end{itemize}
\end{proof}
\begin{note}
    Аналогичное утверждение формулируется для столбцов.
\end{note}
\begin{proposition}
    Строчные и столбцовые ранги не меняются при элементарных преобразованиях.
\end{proposition}
\begin{proof}
    Следует из утверждений 1.3 и 1.4.
\end{proof}
\begin{proposition}
    Пусть $A'$ - подматрица $A$. Тогда $$r_1(A') \leqslant r_1(A)$$ $$r_2(A') \leqslant r_2(A)$$
\end{proposition}
\begin{proof}
    Заметим, что если набор $x$ строк матрицы $A$ - ЛЗ, то и набор из строк любой подматрицы $A'$, содержащийся в строках набора $x$, линейно зависим. Отсюда получаем, что размер набора ЛНЗ строк в $A'$ не больше максимального размера набора ЛНЗ строк в $A$, что равносильно первому неравенству.
    
    Второе неравенство получается применением первого к матрицам $A^T$ и $(A')^T$.
\end{proof}
\begin{theorem}
    Любой набор из $n+1$ столбца высоты n является линейно зависимым.
\end{theorem}
\begin{proof}
    Рассмотрим два случая:
    \begin{itemize}
        \item Если $a_1, \dots, a_n$ - ЛЗ $\Rightarrow$ $a_1, \dots, a_{n+1}$ тоже ЛЗ.
        \item Если $a_1, \dots, a_n$ - ЛНЗ $\Rightarrow$ $A_{n \times n}$ - невырожденная $\Rightarrow a_{n+1}$ - линейная комбинация столбцов $A \Rightarrow$ набор - ЛЗ.  
    \end{itemize}
\end{proof}
\begin{theorem} [о ранге матрицы]
\label{Теорема о ранге матрицы.}
    Для любой матрицы $A_{m \times n}$ выполнено $$r_1 (A) = r_2 (A) = rg(A)$$
\end{theorem}
\begin{proof}
Зафиксируем базисные строки матрицы $A$, а остальные обнулим элементарными операциями. Аналогично поступим со столбцами.
\[
    r_1(A)\{
    \left( 
        \begin{array}{ccc}
         &  &  \\
         & \text{\Large \textit{A}} &  \\
         &  & 
        \end{array}
    \right)
    \longrightarrow
    r_1(A)\{
    \left( 
        \begin{array}{ccc}
         0 & \dots & 0 \\
         & \dots &  \\
         0 & \dots & 0
        \end{array}
    \right)
\]
Тогда на пересечении базисных строк и столбцов мы получим невырожденную подматрицу $A'$, причем $r_1(A) = r_1(A')$ и $r_2(A) = r_2(A')$. Предположим, что $r_1(A') < r_2(A')$, тогда в силу теоремы 2.2 столбцы матрицы $A'$ ЛЗ, противоречие с не вырожденностью. Случай, если $r_1(A') > r_2(A')$ является аналогичным, так как при транспонировании линейные зависимости столбцов и строк сохраняются. Следовательно $r_1(A') = r_2(A')$, а значит и $r_1(A) = r_2(A)$.
\end{proof}
\begin{note}
    Отдельно можно зафиксировать следующий результат: если $rg(A) = r$, то существует невырожденная подматрица размера  $r \times r$.
\end{note}
\begin{definition}
    Минорный ранг - максимальный порядок невырожденной подматрицы.
\end{definition}